% TESAURO 章节写作模板
% 用法：复制到 notes/<chapter>/<chapter>.tex，替换【...】占位后撰写。
% 编译：xelatex <chapter>.tex（需 ctex）
% 约定：先讲故事 -> 形式化 -> 机制/推导 -> 实验图（如有）-> 体系位置 -> 参考资料
%
% 占位请用【...】，不要用 [...]：
%   1) \item [文字] 会被当成枚举标签的可选参数，编号错乱；
%   2) \section*{[文字]} 同理，[ ] 可能被吃掉；
%   3) 正文里不要写 %，否则会截断本行（含右括号）。

\documentclass{article}

\usepackage[UTF8]{ctex}
\usepackage{amsmath,amssymb}
\usepackage{geometry}
\usepackage{graphicx}

\geometry{a4paper, margin=2.5cm}

\begin{document}

% ============================================================
\section*{【中文标题】（【English / Acronym】）}

\subsection*{先讲一个故事}

用一个具体场景建立直觉，不要一上来堆公式。

\textbf{错误做法 / 朴素做法}：【简述】。

\textbf{本章方法的做法}：【简述】。

\textbf{本章的核心问题}：【一句话核心问题】

\[
\boxed{\text{【核心问题或核心目标公式】}}
\]

\subsection*{数学形式化}

定义符号、状态/动作/目标；每个符号用 itemize 解释一次。

\[
\boxed{\text{【形式化定义】}}
\]

\begin{itemize}
  \item $\langle\cdot\rangle$ —— 【含义】；
  \item $\langle\cdot\rangle$ —— 【含义】。
\end{itemize}

% ============================================================
\section*{【第一个技术问题 / 动机】}

\subsection*{【小节标题】}

说明为什么需要这一步；必要时回顾上一章公式。

\[
\text{【关键公式】}
\]

% 若有 notebook 产出的图，放在对应机制段落后：
% \begin{figure}[!htbp]
%   \centering
%   \includegraphics[width=0.88\textwidth]{fig1_xxx.pdf}
%   \caption{【实验设定 + 读图要点（环境、seeds、阴影含义）】。}
%   \label{fig:id}
% \end{figure}

% ============================================================
\section*{【核心机制 / 算法主体】}

\subsection*{机制一：【名称】}

直觉 -> 公式 -> 为什么这样设计。关键公式用 \verb|\boxed{}|。

\[
\boxed{\text{【机制一核心公式】}}
\]

\subsection*{机制二：【名称】}

【机制二正文】

\subsection*{机制三：【名称（可选）】}

【机制三正文】

% ============================================================
\section*{【完整算法 / 流程（可选）】}

\subsection*{【优化目标 / 损失】}

\[
\boxed{\text{【损失或目标】}}
\]

\subsection*{算法流程}

\begin{enumerate}
  \item \textbf{初始化}：【...】；
  \item \textbf{收集数据}：【...】；
  \item \textbf{更新}：【...】。
\end{enumerate}

% ============================================================
\section*{【与相邻方法对比（可选）】}

\subsection*{【本章方法】 vs 【对照方法】}

\begin{itemize}
  \item 【维度一】；
  \item 【维度二】。
\end{itemize}

% ============================================================
\section*{本章在 RL 体系中的位置}

\begin{enumerate}
  \item 【上游：从哪条线走来】；
  \item \underline{【本章】}（本章：【一句话定位】）；
  \item 【下游：通向何处】。
\end{enumerate}

收束段：读者学完应带走什么，以及下一章预告。

\subsection*{参考资料}

\begin{enumerate}
  \item Author et al.\ (Year). Title. \textit{Venue}.
  \item Sutton \& Barto (2018). \textit{Reinforcement Learning: An Introduction} (2nd ed.). Chapter~N.
\end{enumerate}

\end{document}
